\documentclass[11pt]{article}
\input{style-pc}
\usepackage{array}

\begin{document}

\entetesujet{Sujet bac 2015 -- Physique-Chimie -- Série C}

% ============================ CHIMIE ============================
\matiere{CHIMIE}{8 points}

\partie{A : vérification des connaissances}

\noindent\textbf{1. Texte à trous.} Complète la phrase ci-après par les mots : \textit{inférieur, atome, énergie, niveau}.

\textit{Lorsque l'électron de l'$\cdots\cdots$ d'hydrogène passe d'un $\cdots\cdots$ supérieur à un niveau $\cdots\cdots$, l'atome émet de l'$\cdots\cdots$}

\medskip
\noindent\textbf{2. Appariement.} Relie chaque élément-question de la colonne A à un élément-réponse de la colonne B. Exemple : $A_3 = B_3$.

\begin{center}\renewcommand{\arraystretch}{1.5}
\begin{tabular}{|p{0.46\linewidth}|p{0.46\linewidth}|}
\hline
\textbf{Colonne A} & \textbf{Colonne B} \\
\hline
A1) $\ch{^{201}_{84}Po \longrightarrow {}^{206}_{82}Pb + {}^{4}_{2}He}$ & B1) $\ch{2\,Al + 6\,H_3O^+ \longrightarrow 2\,Al^{3+} + H_2}$ \\
\hline
A2) Réaction acido-basique & B2) $\ch{NH_3 + H_3O^+ \longrightarrow NH_4^+ + H_2O}$ \\
\hline
A3) Oxydation du métal zinc & B3) $\ch{Zn \longrightarrow Zn^{2+} + 2\,e^-}$ \\
\hline
A4) Réaction d'oxydo-réduction & B4) Réaction nucléaire spontanée \\
\hline
A5) $\ch{^{238}_{92}U \longrightarrow {}^{234}_{90}Th + {}^{4}_{2}He}$ & B5) Pb est le noyau fils \\
\hline
\end{tabular}
\end{center}

\noindent\textbf{3. Questions à choix multiples.} Choisis la bonne réponse. Exemple : $E = e_5$.
\begin{enumerate}[label=\alph*.]
  \item Une solution aqueuse dont le pH est voisin du pKA est une :
  \begin{enumerate}[label=a\textsubscript{\arabic*}.]
    \item solution réductrice ;
    \item solution tampon ;
    \item solution neutre.
  \end{enumerate}
  \item Lorsque l'atome d'hydrogène est à son niveau d'énergie le plus bas, l'atome est :
  \begin{enumerate}[label=b\textsubscript{\arabic*}.]
    \item à l'état fondamental ;
    \item à l'état excité ;
    \item à l'état ionisé.
  \end{enumerate}
  \item La radioactivité $\beta^-$ correspond à l'émission :
  \begin{enumerate}[label=c\textsubscript{\arabic*}.]
    \item de protons ;
    \item de noyaux d'hélium ;
    \item d'électrons.
  \end{enumerate}
  \item Une solution aqueuse d'acide chlorhydrique est obtenue par dissolution dans l'eau pure :
  \begin{enumerate}[label=d\textsubscript{\arabic*}.]
    \item du dichlore gazeux ;
    \item de gaz chlorure d'hydrogène ;
    \item du chlorure d'aluminium solide.
  \end{enumerate}
\end{enumerate}

\partie{B : application des connaissances (solutions aqueuses)}
Une solution d'éthanamine ($\ch{C_2H_5NH_2}$) de concentration molaire volumique $C_0 = 0{,}1$ mol/L a un pH $= 11{,}8$.

\begin{enumerate}
  \item Vérifie si l'éthanamine est une base forte ou une base faible.
  \item Écris l'équation de la réaction de l'éthanamine avec l'eau.
  \item \begin{enumerate}[label=\alph*.]
    \item Recense les espèces chimiques présentes dans la solution.
    \item Détermine leurs concentrations molaires volumiques.
  \end{enumerate}
  \item Sachant que le couple ion éthanammonium/éthanamine est $\ch{C_2H_5NH_3^+/C_2H_5NH_2}$, calcule le pKa de ce couple acide/base.
\end{enumerate}

% ============================ PHYSIQUE ============================
\matiere{PHYSIQUE}{12 points}

\partie{A : vérification des connaissances}

\noindent\textbf{1. Réarrangement.} La phrase suivante est écrite en désordre. Ordonne-la.

\textit{la trajectoire / géostationnaire / d'un satellite / dans le plan équatorial / est toujours / de la terre.}

\medskip
\noindent\textbf{2. Questions à alternative vrai ou faux.}
\begin{enumerate}[label=\alph*.]
  \item Les ondes mécaniques se propagent dans le vide.
  \item La loi d'Ohm en courant alternatif s'écrit : $U = RI$.
  \item Pour une corde qui est le siège d'ondes stationnaires, l'élongation à l'instant $t$ d'un point vibrant est : $y = 2a\sin\!\left(\dfrac{2\pi x}{\lambda}\right)\cos\!\left(\dfrac{2\pi t}{T}\right)$.
  \item L'équation différentielle du mouvement d'un pendule de torsion est $\ddot{\theta} + \dfrac{C}{J}\theta = 0$.
\end{enumerate}

\partie{B : application des connaissances}
On dispose d'une fourche munie de deux pointes $S_1$ et $S_2$ qui frappent la surface libre d'un liquide au repos. La fourche est liée à un vibreur qui impose deux vibrations sinusoïdales à $S_1$ et $S_2$ en phase, de même amplitude $a = 2\cdot10^{-3}$ m et de même fréquence $N = 100$ Hz. Soit $y_{S_1}(t) = y_{S_2}(t) = a\sin\omega t$, les équations du mouvement des deux sources.

\begin{enumerate}
  \item Établis l'équation du mouvement résultant d'un point $M$ situé à une distance $d_1$ de $S_1$ et $d_2$ de $S_2$.
  \item Calcule le nombre de points d'amplitude maximale qui se forment sur le segment $S_1 S_2$.
  \item Détermine l'état vibratoire d'un point $P$ situé à $d_1 = 3{,}15$ cm de $S_1$ et $d_2 = 4{,}35$ cm de $S_2$.
\end{enumerate}
\noindent On donne : célérité des ondes à la surface du liquide $v = 0{,}6$ m/s ; $S_1 S_2 = 3$ cm.

\partie{C : résolution d'un problème}

\begin{minipage}{0.55\linewidth}
Afin d'évaluer l'impact de la force de frottement sur la vitesse d'un solide, on réalise deux études comparatives en utilisant le dispositif ci-après. Le solide $(S)$, assimilable à un point matériel de masse $m = 10$ g, glisse à l'intérieur de la demi-sphère de centre $O$ et de rayon $r = 1{,}25$ m. On le lâche du point $A$ sans vitesse initiale. Sa position à l'intérieur de la demi-sphère est repérée par $\theta$ (figure ci-contre).
\end{minipage}\hfill
\begin{minipage}{0.42\linewidth}
\centering
\begin{tikzpicture}[>=Stealth,scale=1]
  \def\r{2.4}
  \coordinate (O) at (0,0);
  \coordinate (Bg) at (-\r,0);
  \coordinate (Bb) at (0,-\r);
  \coordinate (S) at ({\r*sin(-215+90)},{-\r*cos(215-180)});
  % S à angle theta de la verticale, côté gauche
  \coordinate (S) at ({-\r*sin(40)},{-\r*cos(40)});
  % diametre
  \draw[dashed] (-\r,0)--(\r,0);
  % demi-cercle bas
  \draw[thick] (-\r,0) arc[start angle=180,end angle=360,radius=\r];
  % rayon vers bas
  \draw[dashed] (O)--(Bb);
  % rayon vers S
  \draw[dashed] (O)--(S);
  % angle theta
  \draw (0,-0.8) arc[start angle=270,end angle={270-40},radius=0.8];
  \node at (-0.35,-0.95){$\theta$};
  % right angle
  \draw (0.25,0)--(0.25,-0.25)--(0,-0.25);
  % r label
  \node[right] at (0.1,-1.3){$r$};
  % points
  \fill (O) circle (1pt) node[above right]{$O$};
  \fill (Bg) circle (1pt) node[left]{$B$};
  \fill (Bb) circle (1pt) node[below]{$B$};
  \fill (S) circle (1.6pt) node[left]{$(S)$};
\end{tikzpicture}
\end{minipage}

\begin{enumerate}
  \item On admet que le solide $(S)$ glisse sans frottement.
  \begin{enumerate}[label=\alph*.]
    \item Exprime sa vitesse au point $M$ en fonction de $g$, $r$ et de $\theta$. Calcule sa valeur numérique $v_B$ au point $B$ ($g = 10$ m$\cdot$s$^{-2}$).
    \item Exprime l'intensité de la réaction $R$ exercée par la demi-sphère sur le solide en fonction de $g$, $r$ et de $\theta$.
    \item Calcule $R$ en $B$.
  \end{enumerate}
  \item En réalité, le solide est soumis à une force de frottement $\vec{f}$ de même direction et de sens opposé au vecteur vitesse $\vec{v}$ du solide. L'intensité de $\vec{f}$ est égale à $1{,}21\cdot10^{-2}$ N.
  \begin{enumerate}[label=\alph*.]
    \item Calcule la vitesse $v'_B$ au point $B$.
    \item Compare $v_B$ et $v'_B$ puis conclus.
  \end{enumerate}
\end{enumerate}

\end{document}
